% !TeX program = lualatex
\documentclass[11pt]{article}
\usepackage[margin=25mm]{geometry}
\usepackage{fontspec}
\setmainfont{Latin Modern Roman}
\usepackage{amsmath,amssymb,mathtools}
\usepackage{unicode-math}
\setmathfont{Latin Modern Math}
\usepackage{xurl,hyperref}
\hypersetup{colorlinks=true,urlcolor=blue}
\setcounter{secnumdepth}{0}
\setlength{\parindent}{0pt}
\setlength{\parskip}{5pt}
\setlength{\emergencystretch}{3em}
\newcommand{\R}{\mathbb R}
\newcommand{\N}{\mathbb N}
\newcommand{\Z}{\mathbb Z}
\newcommand{\Q}{\mathbb Q}
\newcommand{\C}{\mathbb C}
\newcommand{\D}{\mathbb D}
\newcommand{\F}{\mathbb F}
\newcommand{\sref}[2]{\hyperlink{source:#1:#2}{[S#2]}}
\newcommand{\cataloguescope}[1]{#1}
\begin{document}
\section*{The optimal one-third stability exponent for compact-target Brenier maps}\label{statement}
\subsection{Definitions and mathematical statement}
\paragraph{Statement recorded in the supplied source.}
Fix $d\geq 2$, a convex body $K\subset R^d$ with nonempty interior, its uniform probability rho, and a compact target set Y. Is there C(K,Y) such that $||T_\mu -T_\nu ||_(L2(\rho ))\leq C$ $W_2(\mu ,\nu )^{1/3}$ for every pair of target probability measures on Y, including atomic and singular measures? Is no larger universal exponent possible even for cube source and target and exactly three target atoms?
\subsection{Short English statement}
Fix $d\geq 2$, a convex body $K\subset R^d$ with nonempty interior, its uniform probability rho, and a compact target set Y. Is there C(K,Y) such that $||T_\mu -T_\nu ||_(L2(\rho ))\leq C$ $W_2(\mu ,\nu )^{1/3}$ for every pair of target probability measures on Y, including atomic and singular measures? Is no larger universal exponent possible even for cube source and target and exactly three target atoms?
\subsection{Sources}
\begin{itemize}
\item[S1] Original question source or mathematical context.
\url{https://github.com/openai/math/blob/adc7f1241b42e322a6451854ab7e4b4c146bf78a/preprints/Sharp-One-Third-Stability-of-Brenier-Maps-September-25-2026/build/source/sections/introduction.tex}.
\item[S2] OpenAI, manuscript citation page. Source of the formulation motivating this local problem.
\url{https://github.com/openai/math/blob/adc7f1241b42e322a6451854ab7e4b4c146bf78a/preprints/Sharp-One-Third-Stability-of-Brenier-Maps-September-25-2026/README.md}.
\end{itemize}
\end{document}
